%------------------------------------------------------------------------------%
\documentclass[fleqn,utf8,aspectratio=169,12pt]{beamer}

\usepackage{integracao_e_series}

\title{Volumes de Sóidos de Revolução}

%------------------------------------------------------------------------------%
\begin{document}

\addtitle

%------------------------------------------------------------------------------%
\section{Sólidos de Revolução}

%------------------------------------------------------------------------------%
\begin{frame}{Sólidos de Revolução}
  Sólido de revolução é obtido pela rotação de uma
  região plana em torno de um eixo.

  \pause
  \vspace{\baselineskip}
  Seção transversal será um \emph{disco} de raio $R(x)$
  \[
    A(x) = \pi R(x)^2
  \]

  \pause
  \vspace{0.5\baselineskip}
  \emph{Anel} circular de raio externo $R(x)$ e raio interno $r(x)$
  \[
    A(x) = \pi \left(\, R(x)^2 - r(x)^2 \,\right)
  \]
\end{frame}
%------------------------------------------------------------------------------%

\framedgraphic{Sólidos de Revolução}{solido_revolucao_secoes_transversais}{}

%------------------------------------------------------------------------------%
\section{Exemplos}

%------------------------------------------------------------------------------%
\addtocounter{exemplo}{1}
\subsection{Exemplo \theexemplo}

%------------------------------------------------------------------------------%
\begin{frame}{Exemplo \theexemplo}
  Encontre o volume do sólido obtido pela rotação em torno do eixo $x$ \\
  da região sob a curva \(f(x) = \sqrt{x}\) de 0 até 1
\end{frame}

\framedgraphic{Exemplo \theexemplo}{volume_revolucao_sqrt}{}

%------------------------------------------------------------------------------%
\begin{frame}{Exemplo \theexemplo{} -- Área da seção transversal}
  A seção transversal é um círculo de raio
  \[
    R(x)
    \;=\; \sqrt{x}
  \]
  \pause
  com área
  \[
    A(x)
    \;=\; \pi R(x)^2
    \;=\; \pi \left(\sqrt{x}\right)^2
    \;=\; \pi x
  \]
\end{frame}

%------------------------------------------------------------------------------%
\begin{frame}{Exemplo \theexemplo{} -- Volume}
  \[
    V
    \;=\; \int_0^1 \pi x \, dx
    \pause
    \;=\; \pi \left( \frac{x^2}{2}\right) \evalat{0}{1}
    \pause
    \;=\; \frac{\pi}{2}
  \]
\end{frame}

%------------------------------------------------------------------------------%
\addtocounter{exemplo}{1}
\subsection{Exemplo \theexemplo}

%------------------------------------------------------------------------------%
\begin{frame}{Exemplo \theexemplo{}}
  Encontre o volume do sólido obtido pela rotação da região \\
  limitada por $y = x^3$, $y = 8$, e $x = 0$ em torno do eixo $y$
\end{frame}

\framedgraphic{Exemplo \theexemplo{}}{volume_revolucao_y}{}

%------------------------------------------------------------------------------%
\begin{frame}{Exemplo 3 -- Área da seção transversal}
  Seção transversal é um círculo de raio $x$
  \[
    x = \sqrt[3]{y}
  \]

  \pause
  \vspace{\baselineskip}
  Área da seção transversal
  \[
    A(y)
    = \pi R(y)^2
    = \pi \left( \,\sqrt[3]{y}\, \right)^2
    = \pi y^{\sfrac{2}{3}}
  \]
\end{frame}

%------------------------------------------------------------------------------%
\begin{frame}{Exemplo \theexemplo{} -- Volume}
  \[
    V
    \;=\; \int_0^8 \pi y^{\sfrac{2}{3}}dy
    \pause
    \;=\; \pi \left( \frac{3}{5} y^{\sfrac{5}{3}}\right) \evalat{0}{8}
    \pause
    \;=\; \pi \frac{3}{5} 8^{\sfrac{5}{3}}
    \pause
    \;=\; \frac{96 \pi}{5}
  \]
\end{frame}

%------------------------------------------------------------------------------%
\addtocounter{exemplo}{1}
\subsection{Exemplo \theexemplo}

%------------------------------------------------------------------------------%
\begin{frame}{Exemplo \theexemplo{}}
A região $\mathcal{R}$, limitada pelas curvas delimitada pelas curvas
\[
  y = x
  \qquad\text{e}\qquad
  y = x^2
\]
é girada em torno do eixo $x$.

\vspace{\baselineskip}
Encontre o volume do sólido resultante.
\end{frame}

\framedgraphic{Exemplo \theexemplo{}}{solido_revolucao_secoes_transversais}{}

%------------------------------------------------------------------------------%
\begin{frame}{Exemplo 5 -- Área da seção transversal}
As curvas $y = x$ e $y = x^2$ se intersectam nos pontos $(0, 0)$ e $(1, 1)$

\pause
\vspace{0.5\baselineskip}
A seção transversal é um anel

\pause
\vspace{0.5\baselineskip}
Raio interno
\[
  r(x) = x^2
\]

\pause
Raio externo
\[
  R(x) = x
\]

\pause
Área da seção transversal
\[
  A(x)
  \;=\; \pi \left( x^2 - \left(x^2\right)^2 \right)
  \pause
  \;=\; \pi \left( x^2 - x^4 \right)
\]
\end{frame}

%------------------------------------------------------------------------------%
\begin{frame}{Exemplo \theexemplo{} -- Volume}
  \[
    V
    \;=\; \int_0^1 \pi (x^2-x^4) dx
    \pause
    \;=\; \pi \left( \frac{x^3}{3}-\frac{x^5}{5} \right) \evalat{0}{1}
    \pause
    \;=\; \pi \left( \frac{1}{3} - \frac{1}{5} \right)
    \pause
    \;=\; \frac{2\pi}{15}
  \]
\end{frame}

%------------------------------------------------------------------------------%
\addtocounter{exemplo}{1}
\subsection{Exemplo \theexemplo}

%------------------------------------------------------------------------------%
\begin{frame}{Exemplo \theexemplo{}}
  Encontre o volume do sólido obtido pela rotação da região\\
  do exemplo anterior em torno da reta $x=-1$
\end{frame}

\framedgraphic{Exemplo \theexemplo{}}{volume_revolucao_x-1}{}

%------------------------------------------------------------------------------%
\begin{frame}{Exemplo \theexemplo{} -- Área da seção transversal}
  \begin{columns}
  \column{0.4\linewidth}
    A seção transversal é um anel

    \pause
    \vspace{\baselineskip}
    Raio interno
    \[
      r(y) = 1 + y
    \]

    \pause
    \vspace{\baselineskip}
    Raio externo
    \[
      R(y) = 1 + \sqrt{y}
    \]

  \column{0.4\linewidth}
    \pause
    Área da seção transversal
    \begin{align*}
      \onslide<+->{A(y) & = \pi \left(R(y)^2 - r(y)^2\right)                                               \\[2mm]}
      \onslide<+->{     & = \pi \left[\left( 1 + \sqrt{y} \right)^2 - \left( 1 + y \right)^2\right]        \\[2mm]}
      \onslide<+->{     & = \pi \left[\left( 1 + 2\sqrt{y} + y \right) - \left( 1 + 2y + y^2\right)\right] \\[2mm]}
      \onslide<+->{     & = \pi \left( 2\sqrt{y} - y - y^2 \right)}
    \end{align*}
  \end{columns}
\end{frame}

%------------------------------------------------------------------------------%
\begin{frame}{Exemplo \theexemplo{} -- Volume}
\begin{align*}
  \onslide<+->{V & = \int_0^1 A(y) \,dy \\}
  \onslide<+->{  & = \int_0^1 \pi \left( 2\sqrt{y} - y - y^2 \right) dy \\}
  \onslide<+->{  & = \pi \left(
                      \frac{4y^{\sfrac{3}{2}}}{3}
                      - \frac{y^2}{2}
                      - \frac{y^3}{3}
                      \right) \evalat{0}{1} \\}
  \onslide<+->{  & = \pi \left(\frac{4}{3} - \frac{1}{2} - \frac{1}{3}\right) \\}
  \onslide<+->{  & = \frac{\pi}{2}}
\end{align*}
\end{frame}

%------------------------------------------------------------------------------%
\addtocounter{exemplo}{1}
\subsection{Exemplo \theexemplo}

%------------------------------------------------------------------------------%
\begin{frame}{Exemplo \theexemplo{}}
  Encontre o volume do sólido obtido pela rotação da região\\
  do exemplo anterior em torno da reta $y=2$
\end{frame}

\framedgraphic{Exemplo \theexemplo{}}{volume_revolucao_y2}{}

%------------------------------------------------------------------------------%
\begin{frame}{Exemplo \theexemplo{} -- Área da seção transversal}
  A seção transversal é um anel

  \pause
  \vspace{0.5\baselineskip}
  Raio interno
  \[
    r(x) = 2-x
  \]

  \pause
  Raio externo
  \[
    R(x) = 2-x^2
  \]

  \pause
  Área
  \[
    A(x)
    = \pi \left( 2 - x^2 \right)^2 - \pi (2-x)^2
    \pause
    = \pi \left( x^4-5x^2+4x \right)
  \]
\end{frame}

%------------------------------------------------------------------------------%
\begin{frame}{Exemplo \theexemplo{} -- Volume}
  \begin{align*}
    \onslide<+->{V & = \int_0^1 A(x)\, dx           \\}
    \onslide<+->{& = \int_0^1 \pi \left( x^4 - 5x^2 + 4x \right) dx           \\}
    \onslide<+->{& = \pi \left(
        \frac{x^5}{5}
        - 5 \frac{x^3}{3}
        + 4 \frac{x^2}{2}
        \right) \evalat{0}{1}                                    \\}
    \onslide<+->{& = \pi \left(\frac{1}{5} - \frac{5}{3} + \frac{4}{2}\right) \\}
    \onslide<+->{& = \frac{8\pi}{15}}
  \end{align*}
\end{frame}

%------------------------------------------------------------------------------%
\section{Lista Mínima}

%------------------------------------------------------------------------------%
\begin{frame}{Lista Mínima}
  Estudar a Seção 3.3 da Apostila

  Exercícios: 1a-f, 2, 3

  \vfill
  Atenção: A prova é baseada no livro, não nas apresentações
\end{frame}

%------------------------------------------------------------------------------%
\end{document}
%------------------------------------------------------------------------------%
